\section{Síntesis y ampliación}

Esta clase estableció el marco general del aprendizaje por refuerzo profundo. Los puntos clave incluyen:

\begin{enumerate}
    \item El RL profundo estudia los problemas de decisión secuencial y sus soluciones, con una diferencia esencial respecto del aprendizaje supervisado.
    \item Las aplicaciones del RL ya son muy amplias: desde la robótica hasta los videojuegos y el post-training de los LLM.
    \item Marco de modelado básico: estado, acción, política, recompensa, trayectoria.
    \item El objetivo es maximizar la recompensa acumulada esperada $\mathbb{E}_{\tau \sim p_\theta(\tau)} \left[\sum_t r(s_t, a_t)\right]$.
    \item Las siguientes clases presentarán, en orden, imitation learning, policy gradients, actor-critic, Q-learning, offline RL, reward learning y RLHF.
\end{enumerate}

\subsection{Lecturas adicionales}
\begin{itemize}
    \item Sutton \& Barto, \textit{Reinforcement Learning: An Introduction} (2nd edition): libro de texto clásico de RL.
    \item CS234 (Stanford): un curso de RL más orientado a la teoría, complementario a CS224R.
    \item Sergey Levine, \textit{CS285 (Berkeley)}: otro excelente curso de RL profundo.
\end{itemize}

%% --- Fin del contenido --- %%

\end{document}
